\scanpage{70}\mspage{57}
% archive sheet number: 69
% Paragraph continues from the preceding sheet.
du foncteur en sens inverse. $u_W^*$, induit
par $u^*$. Notons que c'est le cas si on
\deleted{suppose que} $u^*$ transforme $M(I'^{\unsure{o|\circ}})_c$ en
$M(I^o)_c$, et quand on suppose que
$I'$ \inserted{est également artinien, et} satisfait \unsure{également} à une des hypothèses
a) b) du th. 3.1. Car alors $(M(I'^o),W(I'^o),C_{I'^o})$
est une catégorie à cofibrations par le \unsure{th.|Th.}
3.1, on peut \unsure{leur|lui} appliquer le lemme
précédent, donc $u_W^*$ et $u_!^W$ se \unsure{réduisent|déduisent} à
partir d'un couple de foncteurs

\diagram{\begin{tikzcd}[ampersand replacement=\&,column sep=large]
M(I^o)_c \arrow[r,dash,shift left=1.4ex,"u_!^c", "\unsure{\text{---}|\to}"{description,near end,above=1pt,font=\tiny}] \& M(I'^o)_c\,, \arrow[l,shift left=1.4ex,"u_c^*"]
\end{tikzcd}}

lesquels sont visiblement adjoints l'un
de l'autre (puisque induits par un
couple de foncteurs adj. $u_!,u^*$), et
de plus compatibles aux localiseurs.
Donc on peut appliquer le lemme des
foncteurs \unsure{adjoint|adjoints} pour les localisés. Pour
résumer :

\uline{Corollaire} 3.5. Supposons que $I,I'$ sont
\deleted{\unsure{loc}} artiniens, et qu'on soit dans l'un
des deux cas \unsure{suivant|suivants}

a) \deleted{$I,I'$ sont \unsure{no}} les espaces $\mathrm{Esp}(I),\mathrm{Esp}(I')$
sont noethériens (p. ex. $I,I'$ finis).
